% ===============================
% Worksheet / Manual Template
% ===============================
\documentclass[12pt, a4paper]{article}

% ===============================
% Metadata
% ===============================
\newcommand*{\CourseCode}{MAT321}
\newcommand*{\CourseTitle}{Real Analysis II}
\newcommand*{\Chapter}{Separable Spaces}
\newcommand*{\Topics}{Dense Set and Perfect Set}
\newcommand*{\DocDate}{March 01, 2026}

\include{header}

% ==========================================
% Document
% ==========================================
\begin{document}
\FrontPage
\Large


\begin{heading}
    Dense Sets
\end{heading}

\vspace{10pt}

\begin{definition}[Dense Set]
A subset $A$ of a metric space $(X,d)$ is said to be \emph{dense} (or \emph{everywhere dense})
in $X$ if the closure of $A$ is $X$, that is,
\[
\overline{A}=X.
\]
\end{definition}

\clearpage

\begin{heading}
    Nowhere Dense Sets
\end{heading}

\vspace{10pt}

\begin{definition}[Nowhere Dense Set]
Let $A$ be a subset of a metric space $(X,d)$. The set $A$ is said to be
\emph{nowhere dense} in $X$ if the interior of its closure is empty, that is,
\[
\operatorname{int}(\overline{A})=\emptyset.
\]
\end{definition}

\vspace{10pt}

\begin{problem}[Basic Properties of Nowhere Dense Sets]
Show that:
\begin{enumerate}
    \item every finite subset of $X$ is nowhere dense;
    \item every subset of a nowhere dense set is nowhere dense;
\end{enumerate}
\end{problem}

\clearpage

\begin{heading}
    Somewhere Dense Sets and Dense-in-itself Sets
\end{heading}

\vspace{10pt}

\begin{definition}[Somewhere Dense Set]
A subset $A$ of $X$ is said to be \emph{somewhere dense} in $X$ if it is not nowhere dense.
Equivalently, $A$ is somewhere dense in $X$ if and only if $\overline{A}$ contains a non-empty
open sphere.
\end{definition}

\vspace{10pt}

\begin{definition}[Dense-in-itself Set]
A subset $A$ of a metric space $(X,d)$ is said to be \emph{dense-in-itself} if every point of
$A$ is a limit point of $A$, that is,
\[
A\subseteq A'.
\]
\end{definition}

\clearpage

\begin{heading}
    Perfect Sets
\end{heading}

\vspace{10pt}

\begin{definition}[Perfect Set]
A subset $A$ of a metric space $(X,d)$ is said to be \emph{perfect} if it is closed and dense-in-itself, that is,
\[A=\overline{A}\quad\text{and}\quad A\subseteq A'.\]
\end{definition}

\vspace{10pt}

\begin{definition}[Other Definition of Perfect Set]
A subset $A$ of a metric space $(X,d)$ is said to be \emph{perfect} if it is closed and has no isolated points.
\end{definition}

\vspace{10pt}

\begin{problem}[Examples of Perfect Sets]
Show that every closed interval, the empty set $\emptyset$, and the whole space $X$ are
perfect sets.
\end{problem}

\clearpage

\begin{definition}[Cantor Set]
The Cantor set is a subset of the real line constructed by iteratively removing the middle third of each interval in the unit interval $[0,1]$.
\end{definition}

\clearpage

\begin{problem}[Cantor Set is Perfect]
Show that the Cantor set is a perfect set.
\end{problem}

\clearpage

\begin{heading}
    Separable Metric Space
\end{heading}

\begin{definition}[Separable Metric Space]
A metric space $(X,d)$ is said to be \emph{separable} if there exists a countable subset of $X$
which is dense in $X$.
\end{definition}

\clearpage

\begin{problem}[Euclidean Space is Separable]
Show that the Euclidean space $\mathbb{R}^n$ is separable.
\end{problem}
\textbf{Hint:} Consider the set of all points in $\mathbb{R}^n$ with rational coordinates as a candidate of a countable dense subset of $\mathbb{R}^n$.
\clearpage

\begin{problem}[$\ell^p$ is Separable]
Show that the metric space $\ell^p$ of all sequences $(x_n)$ of real numbers such that $\sum_{n=1}^\infty |x_n|^p<\infty$ with the metric
\[d_p((x_n),(y_n))=\left(\sum_{n=1}^\infty |x_n-y_n|^p\right)^{1/p}\]
is separable for every $1\leq p<\infty$.
\end{problem}
\textbf{Hint:} Consider the set with finitely many non-zero rational entries as a candidate of a countable dense subset of $\ell^p$.

\clearpage

\begin{problem}[$\ell^\infty$ is Not Separable]
Prove that the metric space $\ell^\infty$ of all bounded sequences with the supremum metric
is not separable.
\end{problem}
\textbf{Hint:} Consider the set of all sequences of $0$'s and $1$'s as a candidate of an uncountable subset of $\ell^\infty$ with pairwise distance at least $1$.


\clearpage

\EndPage
\end{document}